% Abhakseite „Das kann ich“ – Zone nur im Gesamt (\mitzone)
%% TODO Ich-kann-Sätze: Platzhalter wie in den Aufgabendateien
\begin{abhakseite}
\ifdefined\mitzone
\abhakgruppe{Kennst du schon}
\abhak{1}{Räumliches Vorstellen}
\abhak{2}{Flächeninhalt von Rechteck, Dreieck, Trapez und Kreis (Grundflächen)}
\abhak{3}{Längen- und Volumeneinheiten, Kubikdezimeter als Liter, Kubikzentimeter als Milliliter, Umrechnen zwischen beiden}
\abhak{4}{Formel nach einer Größe umstellen (V = G · h → h = V : G)}
\abhak{5}{Multiplizieren mit Dezimalzahlen und π, Runden; Quadrieren und Wurzelziehen}
\abhak{6}{Prozentwert berechnen (zehn Prozent von einem Volumen mit Komma)}
\abhak{7}{Flächeninhalt von Rechteck, Dreieck, Trapez und Kreis (Grundflächen) – Fehler finden}
\abhak{8}{Flächeninhalt von Rechteck, Dreieck, Trapez und Kreis (Grundflächen) – selbst rechnen}
\fi
\abhakgruppe{Einheit 1 · Körper erkennen, Netze, Schrägbilder}
\abhak{9}{Körper und Netze}
\abhak{10}{Grund- und Deckfläche, Seitenflächen benennen (auch beim liegenden Prisma)}
\abhak{11}{Schrägbild lesen: verdeckte Kanten, Maße entnehmen}
\abhak{12}{Körper in ein Schrägbild einzeichnen (Kegel im Quader)}
\abhak{13}{Körper und Netze – Fehler finden, Begründen, Anwendung, Darstellungswechsel}
\abhakgruppe{Einheit 2 · Quader und Würfel}
\abhak{14}{Quader}
\abhak{15}{Oberfläche ohne Deckel (Kiste, Aquarium)}
\abhak{16}{Würfelkante aus V (Kubikzahlen)}
\abhak{17}{Einheiten cm³, dm³ = l, m³ (1 l = 1000 cm³)}
\abhak{18}{Aquarium: Füllmenge in Litern, Füllhöhe}
\abhak{19}{Quader – Fehler finden, Begründen, Anwendung, Darstellungswechsel}
\abhakgruppe{Einheit 3 · Prisma}
\abhak{20}{Prisma}
\abhak{21}{Prisma – Fehler finden, Begründen, Anwendung, Darstellungswechsel}
\abhakgruppe{Einheit 4 · Zylinder}
\abhak{22}{Welche Einheit?}
\abhak{23}{Zylinder}
\abhak{24}{Zylinder – Fehler finden, Begründen, Anwendung, Darstellungswechsel}
\abhakgruppe{Einheit 5 · Zusammengesetzte Körper und Anwendungen}
\abhak{25}{Zusammengesetzt}
\abhak{26}{Füllstand in Prozent}
\abhak{27}{Oberfläche eines zusammengesetzten Körpers}
\abhak{28}{Zusammengesetzt – Fehler finden, Begründen, Anwendung, Darstellungswechsel}
\end{abhakseite}
